\section{Galois-character packets and exact failed routes}

\subsection{Multiquadratic branch codes}

Let $K=\Q(t)$, let $D\subset K^\times/K^{\times2}$ be a $k$-dimensional
space of independent geometric squareclasses, and put
\[
L_D=K(\sqrt d:d\in D).
\]
Let $C_D$ be the smooth projective curve with function field $L_D$, and let
$b$ be the total number of geometric branch points.

The branch-parity vector of a rational function has even weight.  Since
$\mathbf P^1_{\overline{\Q}}$ has no nontrivial unramified double cover, the
branch map embeds $D$ into the even-weight code on $b$ coordinates.  Hence
\[
\boxed{k\le b-1.}
\]
Every branch point has inertia group of order two.  Riemann--Hurwitz gives,
for $k\ge2$,
\[
\boxed{g(C_D)=1+2^{k-2}(b-4).}
\]
Therefore
\[
g(C_D)\ge1+2^{k-2}(k-3).
\]
Thirteen independent quadratic characters require at least fourteen branch
points and hence genus at least
\[
1+2^{11}(14-4)=20481.
\]
Low genus is already lost at the fourth independent character.

\subsection{The complete twist packet}

After tensoring with $\Q$, character idempotents give the exact decomposition
\[
E(L_D)\otimes\Q
\cong
\bigoplus_{d\in D}E^d(K)\otimes\Q.
\]
Distinct character summands are orthogonal under the canonical height.  In the
biquadratic case,
\[
\rank E(K(\sqrt{d_1},\sqrt{d_2}))
=
\rank E(K)+\rank E^{d_1}(K)+\rank E^{d_2}(K)+\rank E^{d_1d_2}(K).
\]
The product-twist channel is often omitted when a search records only the two
multisections used to define the cover.

A single character can carry more than one direction.  The repository
contains an exact family where one nontrivial quadratic eigenspace has rank at
least three.  Thus the branch-code theorem bounds the number of low-genus
characters, not the total rank of their Mordell--Weil packets.

\subsection{A long and useful failed route}

The first proposed mechanism started with a rank-$17$ K3 fibration, constructed
thirteen quadratic multisections, and attempted to split all thirteen over one
low-genus base.  It seemed capable of forcing $17+13=30$.

The genus formula proves that the intended low-genus family cannot arise from
thirteen independent squareclasses.  This is the smallest decisive
obstruction: the fourth independent character already forces genus at least
five.

\subsection{What the failure reveals}

The useful surviving data are:
\begin{itemize}[leftmargin=*]
\item the squareclass of every multisection;
\item all sections in the same character channel;
\item every product-twist channel;
\item the exact Shioda Gram matrix within each channel;
\item branch support, auxiliary genus, and global-solubility data.
\end{itemize}
The search object is therefore a complete \emph{Galois-character height
packet}, not a list of individually constructed points.

\subsection{One-function visibility barrier}

Let $f$ be a rational function on an elliptic curve with divisor
\[
(f)=P_1+\cdots+P_N-NO.
\]
Since a principal divisor maps to the identity in $\operatorname{Pic}^0(E)$,
\[
\boxed{P_1+\cdots+P_N=O.}
\]
Consequently, a construction producing exactly thirty rational points as the
zeros of one function has a forced integral dependence and can span rank at
most twenty-nine.

For a short Weierstrass equation $y^2=R(x)$ and a function
$f=A(x)+yB(x)$, the norm is
\[
\operatorname{Norm}(f)=A(x)^2-R(x)B(x)^2.
\]
The corrected one-function target is therefore a split degree-$31$ identity
\[
A(x)^2-R(x)B(x)^2
=c\prod_{i=1}^{31}(x-r_i),
\qquad r_i\in\Q.
\]
The thirty-one points have one unavoidable relation, leaving room for rank
thirty.  This remains an independent exact construction track.
